\chapter{Termo de Abertura do Projeto}

\textcolor{red}{Termo de abertura do projeto / Project Charter. Um documento publicado pelo iniciador ou patrocinador do projeto que autoriza formalmente a existência de um projeto e fornece ao gerente do projeto a autoridade para aplicar os recursos organizacionais nas atividades do projeto.}

\section{Dados do projeto}
\begin{description}
    \item [Nome do Projeto:] XAROPi
    \item [Data de abertura:] 09/04/2026
    \item [Código:] 1-A
    \item [Patrocinador:] Universidade de Brasília
    \item [Responsável:] Prof. Bruno
    \item [Gerente:] Davi do Egito / 231030421 / 231030421@aluno.unb.br
\end{description}

\section{Objetivos}
\textcolor{red}{O que a empresa pretende obter com a realização do projeto. Descrever o que se pretende realizar para resolver o problema central ou explorar a oportunidade identificada. Para a correta definição do objetivo siga a regra "SMART":
\begin{description}
    \item [\textit{Specific} (específico):] Deve ser redigido de forma clara, concisa e compreensiva;
    \item [\textit{Measurable} (mensurável):] O objetivo específico deve ser mensurável, ou seja, 
possível de ser medido por meio de um ou mais indicadores;
    \item [\textit{Agreed} (acordado):] Deve ser acordado com as partes interessadas, ou seja, as áreas envolvidas na empresa: P\&D, Produção, Comercial, Marketing, Financeira, Jurídica, Manutenção, ambiental, entre outras;
    \item [\textit{Realistic} (realista):] Deve estar centrado na realidade, no que é possível de ser feito considerando as premissas e restrições existentes, como: orçamento e tempo;
    \item [\textit{Time Bound} (Limitado no tempo):] Deve ter um prazo determinado para sua finalização.
\end{description}
}

\section{Mercado-alvo}

\textcolor{red}{Pessoas, empresas, instituições etc. que usufruirão dos produtos, serviços e resultados gerados pelo projeto, cujos requisitos (tópico abaixo) devem atender as suas necessidades. Podem ser internas ou externas à organização, mas, merecem destaque especial, pois, o projeto está sendo feito para atendê-los de forma direta ou indireta.}

\section{Requisitos}

\textcolor{red}{Listar os fatos que são essenciais para o consumidor adquirir o produto, exemplo: cor, tamanho, material, tempo de vida-útil, dispositivo de segurança \textbf{x}, entre outros.}


\paragraph*{Instruções do Grupo}
Ao escrever requisitos, atentem-se às seguintes recomendações:
\begin{itemize}
    \item Atomicidade: cada requisito deve ser indivisível. Caso descreva mais de uma funcionalidade (ex.: "o sistema deve fazer A e B"), ele deve ser separado em dois ou mais requisitos independentes;
    \item Mensurabilidade: os requisitos devem ser objetivos e verificáveis. Evitem termos vagos ou subjetivos (como "rápido" ou "eficiente"). Prefiram descrições quantificáveis, como "latência máxima de 600 ms";
    \item Classificação:
    \begin{itemize}
        \item Requisitos funcionais descrevem comportamentos ou funcionalidades do sistema (ex.: autenticação de usuário, exibição de dados);
        \item Requisitos não funcionais definem restrições ou atributos de qualidade (ex.: desempenho, confiabilidade, usabilidade, consumo de energia);
    \end{itemize}
    \item Padronização: todos os requisitos devem seguir a estrutura definida no projeto, utilizando os comandos:
    \begin{itemize}
        \item \texttt{\textbackslash rf\{<título>\}\{<descrição>\}};
        \item \texttt{\textbackslash rnf[<rótulo>]\{<título>\}\{<descrição>\}}.
    \end{itemize}
\end{itemize}


\begin{table}[htpb]
\begin{center}
\caption{Requisitos Funcionais.}
\begin{tabular}{l l p{8cm}}
    \toprule
    ID & Nome & Descrição \\
    \midrule
    \rf[cadastro]{Cadastro de produtos}{O sistema deve permitir cadastro de produtos com nome, SKU e preço unitário.} \\   
    \rf[login]{Autenticação de usuário}{O sistema deve autenticar o usuário via e-mail e senha, com bloqueio após 5 tentativas falhas.} \\
    \bottomrule
\end{tabular}
\end{center}
\end{table}

\begin{table}[htpb]
\begin{center}
\caption{Requisitos Não Funcionais.}
\begin{tabular}{l l p{8cm}}
    \toprule
    ID & Nome & Descrição \\
    \midrule
    \rnf[usabilidade_feedback]{Feedback do sistema}{O sistema deve fornecer feedback visual imediato (em até 200 ms) após cada interação do usuário.} \\
    \bottomrule
\end{tabular}
\end{center}
\end{table}

Exemplo de referência a um requisito: \cref{req:login}.

\begin{longtable}{l l p{8cm}}
    \caption{Requisitos Funcionais -- Hardware} \\
    \toprule
    ID & Nome & Descrição \\
    \midrule
    \endfirsthead
    \endhead
    \endfoot
    \bottomrule
    \endlastfoot

    \rf[tof-paredes]{Detecção de paredes do labirinto}{
        O sistema deve detectar a presença ou ausência de paredes
        nas direções frontal, lateral esquerda e lateral direita
        do micromouse, compatível com as dimensões das células de
        18~cm × 18~cm e paredes de 5~cm de altura.} \\
    \rf[motor-pwm]{Controle de velocidade dos motores}{
        O sistema deve controlar individualmente a velocidade de
        cada motor DC por meio de sinal PWM enviado ao driver
        de motor.} \\
    \rf[motor-direcao]{Controle de direção dos motores}{
        O sistema deve controlar o sentido de rotação de cada motor
        DC de forma independente, permitindo movimentos de avanço,
        recuo e curvas.} \\
    \rf[linha-reta]{Deslocamento em linha reta}{
        O sistema deve mover o micromouse em linha reta dentro dos
        corredores do labirinto, mantendo trajetória centralizada
        entre as paredes.} \\
    \rf[curva]{Execução de curvas}{
        O sistema deve executar curvas de 90° e 180° de forma
        precisa, compatível com as dimensões das células de
        18~cm × 18~cm do labirinto.} \\
    \rf[botao]{Acionamento por botão}{
        O sistema deve iniciar e interromper a navegação por meio
        de um botão físico acessível externamente no micromouse} \\
    \rf[bat-consumo]{Autonomia de bateria}{
        A bateria do sistema deve fornecer energia suficiente para
        completar os três labirintos (4×4, 8×8 e 16×16
        células) em sequência sem necessidade de recarga.} \\
    \rf[bat-alerta]{Alerta de bateria baixa}{
        O sistema deve detectar e sinalizar quando a tensão da
        bateria atingir nível insuficiente para continuidade
        segura da navegação.} \\
    \rf[buzzer]{Sinalização sonora ao atingir objetivo}{
        O sistema deve acionar um buzzer ao detectar que o micromouse
        atingiu a sala objetivo do labirinto.} \\
    \rf[wifi-telemetria]{Transmissão de telemetria}{
        O sistema deve transmitir dados de telemetria ao sistema
        web em tempo real.} \\
\end{longtable}

\section{Justificativa}

\textcolor{red}{Informar o problema ou a oportunidade (necessidade) que justifica o porquê de o projeto ser realizado. Por exemplo: atende uma demanda específica do consumidor final; supre uma necessidade do mercado comercializador; é um diferencial X para o órgão regulamentador.}

\section{Indicadores}

\textcolor{red}{Listar até 10 indicadores que determinam o mercado consumidor do produto desenvolvido: exemplo: 1) n\textdegree de alunos da FGA que utilizam ônibus às 18:00; 2) n\textdegree de usuários do restaurante universitários, 3) número de idosos classificados como público-alvo no DF e no estado de Goiás, 4) n\textdegree de empresas de segurança registradas no DF etc.}
% \section{Aprovações}

% \begin{tabular}{ l l }
%   \textbf{Patrocinador:} & \_\_\_\_\_\_\_\_\_\_\_\_\_\_\_\_\_\_\_\_\_\_\_\_\_\_\_\_\_\_\_\_\_\_ \\
%   & \\
%   \textbf{CEO:} & \_\_\_\_\_\_\_\_\_\_\_\_\_\_\_\_\_\_\_\_\_\_\_\_\_\_\_\_\_\_\_\_\_\_ \\
%   & \\
%   \textbf{Gerente:} & \_\_\_\_\_\_\_\_\_\_\_\_\_\_\_\_\_\_\_\_\_\_\_\_\_\_\_\_\_\_\_\_\_\_ \\
% \end{tabular}
